\begin{abstract}
QSAR models for Ames mutagenicity are increasingly used in regulatory risk assessment under EU REACH, K-REACH, and ICH M7, 
but their reported performance varies widely across evaluation protocols.\cite{Honma_2019,Furuhama_2023,Wu_MoleculeNet_2018} 
In this study, 225 model configurations---seven algorithms, five feature sets, and three preprocessing 
scenarios---were systematically evaluated across three genotoxicity datasets: Ames 
($n = 13{,}560$), \textit{in vitro} chromosome aberration ($n = 1{,}619$), and \textit{in vivo} micronucleus 
($n = 2{,}153$). Using a fixed model configuration (XGBoost, compact features, raw SMILES), 
Ames MCC dropped from 0.670 under random-split CV to 0.624 under scaffold-split CV. 
Cross-source leave-domain-out (LDO) evaluation produced markedly larger declines: 0.234 
(drug$\to$industrial) and 0.122 (industrial$\to$drug). The source heterogeneity gap 
($\Delta = 0.390$) considerably exceeded the split-method gap ($\Delta = 0.046$). A simple 
source-based classifier reached MCC~$= 0.608$ with no structural information, reflecting a 
19.5-fold difference in positive rates between drug-sourced (58.3\%) and industrial-sourced 
(3.0\%) compounds. Threshold decomposition showed that prior shift accounted for only part 
of this gap, with oracle MCC plateauing at 0.27--0.36. Algorithm robustness under source 
shift varied considerably: Random Forest maintained MCC~$= 0.325$ in the hardest direction, 
whereas XGBoost collapsed to 0.122. \textit{In vivo} micronucleus models showed strong ranking 
performance (ROC-AUC~$= 0.860$, PR-AUC~$= 0.237$) but unreliable classification (MCC 
95\% CI crossed zero). These results call for a source-aware evaluation cascade as a 
standard reporting requirement for mutagenicity QSAR.
\end{abstract}
